\chapter{Affine Connections; Riemannian Connections}
\label{chap:dc-affine-connections-riemannian-connections}

Faithful transcription of do Carmo, \emph{Riemannian Geometry}. Each numbered
statement keeps do Carmo's number, kind and (name); proofs keep his explicit
cross-references ("by \cref{def:dc-ch2-2-1}", "from \cref{prop:dc-ch2-2-6}") so the dependency
graph is recoverable from the text.

\section{Introduction}

Historically, the notion of covariant derivative arose from the Levi-Civita
parallelism (1917). For a surface $S\subset\mathbb{R}^3$ and a vector field $V$
along a curve $c$ tangent to $S$, the usual derivative $\frac{dV}{dt}$ need not
lie in $T_{c(t)}S$; the \emph{covariant derivative} $\frac{DV}{dt}$ is its orthogonal
projection onto $T_{c(t)}S$, depending only on the first fundamental form. A
field $V$ is \emph{parallel} if $\frac{DV}{dt}\equiv 0$. The covariant derivative lets
one define geodesics and curvature in settings more general than Riemannian
manifolds, leading to the notion of \emph{affine connection} (\cref{def:dc-ch2-2-1}). Our
interest in affine connections rests in the fact (\cref{thm:dc-ch2-3-6}) that a choice of
Riemannian metric on $M$ uniquely determines a certain affine connection on $M$.

\section{Affine connections}

Throughout, $\mathcal{X}(M)$ denotes the set of all $C^\infty$ vector fields on
$M$ and $\mathcal{D}(M)$ the ring of real-valued $C^\infty$ functions on $M$.

\begin{definition}[affine connection]
  \label{def:dc-ch2-2-1}
  \dcref{ch2:2.1}
  \lean{Riemannian.AffineConnection}
  \leanok
  An \emph{affine connection} $\nabla$ on a differentiable manifold $M$ is a mapping
  
  $$
  \nabla:\mathcal{X}(M)\times\mathcal{X}(M)\to\mathcal{X}(M),
  $$
  
  denoted $(X,Y)\mapsto\nabla_X Y$, satisfying:
  - i) $\nabla_{fX+gY}Z = f\nabla_X Z + g\nabla_Y Z$.
  - ii) $\nabla_X(Y+Z)=\nabla_X Y+\nabla_X Z$.
  - iii) $\nabla_X(fY)=f\nabla_X Y + X(f)Y$,
  
  for all $X,Y,Z\in\mathcal{X}(M)$ and $f,g\in\mathcal{D}(M)$.
\end{definition}

\begin{proposition}[covariant derivative along a curve]
  \label{prop:dc-ch2-2-2}
  \dcref{ch2:2.2}
  \uses{def:dc-ch2-2-1, lem:dc-ch2-2-2-coord, lem:dc-ch2-2-2-c-coord, lem:dc-ch2-2-2-c-bridge}
  \lean{Riemannian.covariantDerivCoord, Riemannian.covariantDerivCoord_add,
    Riemannian.covariantDerivCoord_smul, Riemannian.covariantDerivCoord_induced,
    Riemannian.covariantDerivCoord_unique}
  \leanok
  Let $M$ be a differentiable manifold with an affine connection $\nabla$. There
  exists a unique correspondence which associates to a vector field $V$ along the
  differentiable curve $c:I\to M$ another vector field $\frac{DV}{dt}$ along $c$,
  called the \emph{covariant derivative} of $V$ along $c$, such that:
  - a) $\frac{D}{dt}(V+W)=\frac{DV}{dt}+\frac{DW}{dt}$.
  - b) $\frac{D}{dt}(fV)=\frac{df}{dt}V+f\frac{DV}{dt}$, where $W$ is a vector field along $c$ and $f$ a differentiable function on $I$.
  - c) If $V$ is induced by a vector field $Y\in\mathcal{X}(M)$, i.e. $V(t)=Y(c(t))$, then $\frac{DV}{dt}=\nabla_{dc/dt}Y$.
\end{proposition}
\begin{proof}
  \leanok
  Suppose first a correspondence satisfying (a), (b), (c) exists. In a
  coordinate system $\mathbf{x}:U\subset\mathbb{R}^n\to M$ write $V=\sum_j v^j X_j$
  with $X_j=\frac{\partial}{\partial x_j}(c(t))$. By a) and b),
  $\frac{DV}{dt}=\sum_j \frac{dv^j}{dt}X_j+\sum_j v^j\frac{DX_j}{dt}$. By c) and (i)
  of \cref{def:dc-ch2-2-1}, $\frac{DX_j}{dt}=\nabla_{dc/dt}X_j=\sum_i\frac{dx_i}{dt}\nabla_{X_i}X_j$.
  Therefore
  
  $$
  \frac{DV}{dt}=\sum_j\frac{dv^j}{dt}X_j+\sum_{i,j}\frac{dx_i}{dt}v^j\nabla_{X_i}X_j.\qquad(1)
  $$
  
  Expression (1) shows uniqueness. For existence, define $\frac{DV}{dt}$ in
  $\mathbf{x}(U)$ by (1); it has the desired properties, and on overlaps the local
  definitions agree by uniqueness, so the definition extends over all of $M$.
  $\blacksquare$
\end{proof}

\begin{lemma}[covariant derivative in coordinates; properties (a), (b)]
  \label{lem:dc-ch2-2-2-coord}
  \dcref{ch2:2.2}
  \uses{def:dc-ch2-3-7-christoffel}
  \lean{Riemannian.covariantDerivCoord, Riemannian.covariantDerivCoord_add,
    Riemannian.covariantDerivCoord_smul}
  \leanok
  Fix a Riemannian metric $g$ and a chart at $\alpha$, with Christoffel symbols
  $\Gamma^k_{ij}$ and Christoffel contraction $\Gamma(v,w)(y)=\sum_k(\sum_{ij}\Gamma^k_{ij}(y)v^iw^j)e_k$.
  For a coordinate curve $u:I\to\mathbb{R}^n$ and a coordinate vector field
  $V:I\to\mathbb{R}^n$ along it, the covariant derivative of
  \cref{prop:dc-ch2-2-2} is given by formula (1),
  $$\frac{DV}{dt}=\dot V+\Gamma(\dot u,V)(u).$$
  This coordinate operator is additive in $V$ and satisfies the Leibniz rule
  $\frac{D}{dt}(fV)=\dot f\,V+f\frac{DV}{dt}$ for differentiable $f:I\to\mathbb{R}$,
  i.e. properties (a) and (b) of \cref{prop:dc-ch2-2-2}, for the Levi-Civita
  connection of $g$.
\end{lemma}

\begin{lemma}[property (c) in coordinates: induced fields]
  \label{lem:dc-ch2-2-2-c-coord}
  \dcref{ch2:2.2}
  \uses{lem:dc-ch2-2-2-coord, def:dc-ch2-3-7-christoffel}
  \lean{Riemannian.chartCovariantDeriv, Riemannian.covariantDerivCoord_induced}
  \leanok
  Read the covariant derivative of a coordinate field $Y:E\to E$ in a direction $v$ at a base
  point $y$ as $(\nabla_v Y)(y)=DY(y)\,v+\Gamma(v,Y(y))(y)$ (Remark 2.3,
  $\nabla_X Y=\sum_k(\sum_{ij}x_iy_j\Gamma^k_{ij}+X(y_k))X_k$). When the vector field $V$ along
  the coordinate curve $u$ is \emph{induced} by $Y$, $V(t)=Y(u(t))$, the covariant derivative
  along $u$ equals the coordinate covariant derivative of $Y$ in the direction $\dot u(t)$:
  $$\frac{D}{dt}\big(Y\circ u\big)=(\nabla_{\dot u}Y)(u).$$
  This is do Carmo's key step $\frac{DX_j}{dt}=\nabla_{dc/dt}X_j$ (property (c) at the level of
  the coordinate connection): the chain rule $\frac{d}{dt}Y(u(t))=DY(u(t))\,\dot u(t)$ turns the
  $\dot V$ term into the directional-derivative term $DY\,\dot u$, while the Christoffel
  contraction matches on the nose.
\end{lemma}

\begin{lemma}[formula (10): the Koszul reduction on a coordinate frame]
  \label{lem:dc-ch2-2-2-c-formula10}
  \dcref{ch2:2.2}
  \uses{lem:dc-ch2-3-6-koszul, def:dc-ch2-3-7-christoffel, thm:dc-ch2-3-6}
  \lean{Riemannian.christoffel_bridge_inner, Riemannian.christoffel_bridge_vector}
  \leanok
  Let $\nabla$ be a Levi-Civita connection for $g$ and let $X_1,\dots,X_n$ be smooth vector
  fields that form a coordinate frame at $p$: the values $\{X_a(p)\}$ are the chart-basis
  vectors (hence a basis of $T_pM$), the pairwise Lie brackets vanish, $[X_a,X_b](p)=0$, and the
  directional derivatives of the metric are the partial derivatives of the chart Gram matrix,
  $X_r\langle X_a,X_b\rangle(p)=\partial_r G_{ab}$. Then the inner product of the frame field
  $X_k$ with $\nabla_{X_i}X_j$ is the Gram contraction of the chart Christoffel symbols
  $\Gamma^m_{ij}=$ \texttt{chartChristoffel},
  $$\langle X_k,\nabla_{X_i}X_j\rangle=\sum_m G_{km}\,\Gamma^m_{ij},$$
  and hence
  $$\nabla_{X_i}X_j=\sum_m\Gamma^m_{ij}\,X_m,$$
  which is do Carmo's formula (10).
  \begin{proof}
    \leanok
    Apply the Koszul formula (9) (\cref{lem:dc-ch2-3-6-koszul}) with the covariant derivative
    $\nabla_{X_i}X_j$ tested against $Z=X_k$. The three Lie-bracket terms of (9) vanish because
    $[X_a,X_b](p)=0$, and the three directional derivatives become
    $X_r\langle X_a,X_b\rangle(p)=\partial_r G_{ab}$, so the right-hand side of (9) is exactly the
    metric-derivative combination $\partial_i G_{kj}+\partial_j G_{ki}-\partial_k G_{ij}$
    (after using the symmetry $G_{ab}=G_{ba}$). By the Gram--Christoffel contraction
    (\texttt{chartGram\_christoffel\_contraction}) this equals $2\sum_m G_{km}\Gamma^m_{ij}$,
    which gives the inner form. For the vector form, both $\nabla_{X_i}X_j$ and
    $\sum_m\Gamma^m_{ij}X_m$ have the same inner product against every $X_k$
    (using $G_{km}=\langle X_k,X_m\rangle$), and $\{X_m(p)\}$ is a basis, so by non-degeneracy of
    $g$ the two vectors are equal.
  \end{proof}
\end{lemma}

\begin{lemma}[property (c): chart-Christoffel bridge to the abstract connection]
  \label{lem:dc-ch2-2-2-c-bridge}
  \dcref{ch2:2.2}
  \uses{lem:dc-ch2-2-2-c-coord, lem:dc-ch2-2-2-c-formula10, thm:dc-ch2-3-6, def:dc-ch2-3-7-christoffel}
  \lean{Riemannian.exists_chartFrame_leviCivita_christoffel}
  \leanok
  The remaining ingredient of the manifold-level property (c): the coordinate covariant
  derivative $\nabla_v Y$ of \cref{lem:dc-ch2-2-2-c-coord}, built from the metric's chart
  Christoffel symbols $\Gamma^k_{ij}=$ \texttt{chartChristoffel}, coincides with the abstract
  Levi-Civita connection $\nabla$ of \cref{thm:dc-ch2-3-6} read in the chart, i.e.
  $\Gamma^k_{ij}=\Gamma(\nabla)^k_{ij}$. Concretely: at any chart-source point $p$ there are
  global smooth vector fields $X_1,\dots,X_n$ realising the chart basis at $p$ for which the
  Levi-Civita connection $\nabla=$ \texttt{leviCivitaConnection} satisfies do Carmo's formula (10)
  with the metric's own chart Christoffel symbols, $\nabla_{X_i}X_j=\sum_m\Gamma^m_{ij}X_m$.
  \begin{proof}
    \leanok
    The frame is the concrete chart frame \texttt{chartBasisVec} $\alpha\,\cdot$ extended to
    global smooth vector fields agreeing with it near $p$
    (\texttt{exists\_smoothVectorField\_eventuallyEq}). Its two analytic identities are
    discharged germ-locally. \emph{Bracket vanishing} $[X_a,X_b](p)=0$: near $p$ the frame
    field is the chart-pullback of the constant model field $e_i$ (the differential of the chart
    sends $X_i(p)$ to $e_i$, \texttt{mfderiv\_extChartAt\_chartBasisVecFiber}), so by naturality
    of the Lie bracket (\texttt{mpullback\_mlieBracket}) the frame bracket is the chart-pullback
    of the bracket of two constant fields, which is $0$. \emph{Directional-derivative $=$
    partial-derivative}: the chart Gram entry $q\mapsto G_{ab}(q)$ agrees near $p$ with
    $G_{ab}\circ\varphi_\alpha$, so the chart chain rule and the same differential identity give
    $X_r\langle X_a,X_b\rangle(p)=\partial_r G_{ab}$. Feeding these to
    \cref{lem:dc-ch2-2-2-c-formula10} (\texttt{christoffel\_bridge\_vector}) with
    $\nabla=$ \texttt{leviCivitaConnection} (Levi-Civita by
    \texttt{isLeviCivita\_of\_koszulDual}) yields formula (10), identifying the coordinate
    Christoffel symbols with those of the abstract Levi-Civita connection. This closes the full
    property (c) and, with it, \cref{prop:dc-ch2-2-2}.
  \end{proof}
\end{lemma}

\begin{remark}[Remark 2.3]
  \label{rem:dc-ch2-2-3}
  \dcref{ch2:2.3}
  \uses{rem:dc-ch0-5-7, def:dc-ch2-2-1}
  The last line of (c) makes sense, since $\nabla_X Y(p)$ depends on the value of
  $X(p)$ and the value of $Y$ along a curve tangent to $X$ at $p$. Part (iii) of
  \cref{def:dc-ch2-2-1} shows that the affine connection is a \emph{local} notion (cf. \cref{rem:dc-ch0-5-7}). In coordinates $(x_1,\dots,x_n)$ about $p$ with
  $X=\sum_i x_i X_i$, $Y=\sum_j y_j X_j$, $X_i=\frac{\partial}{\partial x_i}$,
  
  $$
  \nabla_X Y=\sum_{ij}x_i y_j\nabla_{X_i}X_j+\sum_{ij}x_i X_i(y_j)X_j.
  $$
  
  Setting $\nabla_{X_i}X_j=\sum_k\Gamma_{ij}^k X_k$ with $\Gamma_{ij}^k$
  differentiable,
  
  $$
  \nabla_X Y=\sum_k\Big(\sum_{ij}x_i y_j\Gamma_{ij}^k+X(y_k)\Big)X_k,
  $$
  
  which shows $\nabla_X Y(p)$ depends on $x_i(p)$, $y_k(p)$ and the derivatives
  $X(y_k)(p)$.
\end{remark}

\begin{remark}[Remark 2.4]
  \label{rem:dc-ch2-2-4}
  \dcref{ch2:2.4}
  \uses{prop:dc-ch2-2-2}
  \cref{prop:dc-ch2-2-2} shows that the choice of an affine connection on $M$ leads to a
  bona fide (i.e. satisfying (a) and (b)) derivative of vector fields along curves,
  furnishing a manner of differentiating vectors along curves; in particular one
  may speak of the acceleration of a curve in $M$.
\end{remark}

\begin{definition}[parallel vector field]
  \label{def:dc-ch2-2-5}
  \dcref{ch2:2.5}
  \uses{lem:dc-ch2-2-2-coord}
  \lean{Riemannian.IsParallelCoord}
  \leanok
  Let $M$ be a differentiable manifold with an affine connection $\nabla$. A vector
  field $V$ along a curve $c:I\to M$ is called \emph{parallel} when $\frac{DV}{dt}=0$ for
  all $t\in I$. In coordinates (\cref{lem:dc-ch2-2-2-coord}) this is the linear
  system $\dot V=-\Gamma(\dot u,V)(u)$.
\end{definition}

\begin{lemma}[the parallel-transport coefficient is a bounded linear operator]
  \label{lem:dc-ch2-2-6-coeff}
  \dcref{ch2:2.6}
  \uses{def:dc-ch2-2-5}
  \lean{Riemannian.chartChristoffelContractionRight,
    Riemannian.lipschitzWith_neg_chartChristoffelContraction}
  \leanok
  Fix a velocity $v$ and a base point $y$ in the chart at $\alpha$. The map
  $w\mapsto\Gamma(v,w)(y)$ is linear in $w$, hence --- the model space being
  finite-dimensional --- a bounded linear operator $A(v,y)\colon E\to E$.
  Consequently the right-hand side $w\mapsto-\Gamma(\dot u(t),w)(u(t))$ of the
  parallelism system~(2) is globally Lipschitz in $w$ with constant
  $\lVert A(\dot u(t),u(t))\rVert$. This is the a-priori bound that drives
  Grönwall's inequality in the uniqueness of parallel transport, and the estimate
  underpinning the (linear) global existence of the solution of~(2).
\end{lemma}

\begin{proposition}[parallel transport]
  \label{prop:dc-ch2-2-6}
  \dcref{ch2:2.6}
  \uses{def:dc-ch2-2-5, lem:dc-ch2-2-6-exist, lem:dc-ch2-2-6-unique, lem:dc-ch2-2-6-transport}
  \lean{Riemannian.parallelTransport}
  \leanok
  Let $M$ be a differentiable manifold with an affine connection $\nabla$. Let
  $c:I\to M$ be a differentiable curve and let $V_0$ be a vector tangent to $M$ at
  $c(t_0)$, $t_0\in I$ (i.e. $V_0\in T_{c(t_0)}M$). Then there exists a unique
  parallel vector field $V$ along $c$ with $V(t_0)=V_0$; $V(t)$ is called the
  \emph{parallel transport} of $V(t_0)$ along $c$.
\end{proposition}
\begin{proof}
  \leanok
  \uses{lem:dc-ch2-2-6-exist, lem:dc-ch2-2-6-unique, lem:dc-ch2-2-6-transport}
  Suppose first that the theorem is proved when $c(I)$ is contained in a
  local coordinate neighborhood. By compactness, for any $t_1\in I$,
  $c([t_0,t_1])$ can be covered by finitely many coordinate neighborhoods, and
  uniqueness makes the local definitions coincide on overlaps, so it suffices to
  prove the theorem when $c(I)$ lies in one coordinate neighborhood
  $\mathbf{x}(U)$. Writing $V=\sum v^j X_j$, parallelism gives
  $0=\frac{DV}{dt}=\sum_j\frac{dv^j}{dt}X_j+\sum_{i,j}\frac{dx_i}{dt}v^j\nabla_{X_i}X_j$.
  Putting $\nabla_{X_i}X_j=\sum_k\Gamma_{ij}^k X_k$ yields the system of $n$ linear
  ODEs in $v^k(t)$,
  
  $$
  0=\frac{dv^k}{dt}+\sum_{i,j}\Gamma_{ij}^k v^j\frac{dx_i}{dt},\quad k=1,\dots,n,\qquad(2)
  $$
  
  which possesses a unique solution with $v^k(t_0)=v_0^k$. Since the system is
  linear, the solution is defined for all $t\in I$, proving existence and
  uniqueness. $\blacksquare$
\end{proof}

\begin{lemma}[existence of parallel transport]
  \label{lem:dc-ch2-2-6-exist}
  \dcref{ch2:2.6}
  \uses{def:dc-ch2-2-5, lem:dc-ch2-2-6-coeff}
  \lean{Riemannian.exists_isParallelCoord_Icc}
  \leanok
  On a compact time interval $[a,b]$, for any prescribed initial vector $V_0$ there is a
  coordinate vector field $V$ along the coordinate curve $u$ with $V(a)=V_0$ that is
  parallel, i.e. solves the first-order linear system $\dot V=-\Gamma(\dot u,V)(u)$ on
  $[a,b]$. This is the existence half of \cref{prop:dc-ch2-2-6}. The parallelism system is
  the linear ODE $\dot V=A(t)V$ with coefficient
  $A(t)=-\Gamma(\dot u(t),\,\cdot\,)(u(t))$, a continuous bounded operator by
  \cref{lem:dc-ch2-2-6-coeff}; its global solution is obtained by solving Picard--Lindelöf on
  short subintervals and gluing across a partition of $[a,b]$ (this global linear-ODE
  existence, \texttt{Riemannian.LinearODE.exists\_hasDerivWithinAt\_Icc}, is not in mathlib's
  purely local Picard--Lindelöf theory). The only hypotheses on the curve are that the
  coefficient is continuous and bounded on $[a,b]$; both hold for a $C^1$ curve with
  continuous Christoffel symbols by compactness.
\end{lemma}

\begin{lemma}[uniqueness of parallel transport]
  \label{lem:dc-ch2-2-6-unique}
  \dcref{ch2:2.6}
  \uses{def:dc-ch2-2-5, lem:dc-ch2-2-2-coord, lem:dc-ch2-2-6-coeff}
  \lean{Riemannian.isParallelCoord_eqOn_Icc, Riemannian.isParallelSol_eqOn_Icc}
  \leanok
  On a compact time interval $[a,b]$, two parallel coordinate vector fields along
  the same coordinate curve $u$ that agree at the initial time $a$ agree on all of
  $[a,b]$. This is the uniqueness half of \cref{prop:dc-ch2-2-6}: the parallelism
  equation (2) is a first-order linear ODE whose right-hand side is uniformly
  Lipschitz by \cref{lem:dc-ch2-2-6-coeff}, so Grönwall's inequality forces the two
  solutions to coincide. (Existence — the remaining half — is the Picard--Lindelöf
  solution of the same linear system.)
\end{lemma}

\begin{lemma}[parallel transport is a linear isomorphism]
  \label{lem:dc-ch2-2-6-transport}
  \dcref{ch2:2.6}
  \uses{lem:dc-ch2-2-6-exist, lem:dc-ch2-2-6-unique}
  \lean{Riemannian.parallelTransport, Riemannian.exists_parallelTransport_spec}
  \leanok
  On a compact time interval $[a,b]$ the assignment $P_c:V_0\mapsto V(b)$, sending an
  initial vector $V_0\in T_{c(a)}M$ to the value at $b$ of the parallel field $V$ along
  the coordinate curve $u$ with $V(a)=V_0$, is a well-defined \emph{linear isomorphism}
  $T_{c(a)}M\to T_{c(b)}M$ — do Carmo's \emph{parallel transport}. Well-definedness and
  the value $V(b)=P_c V_0$ realize the parallel field of \cref{lem:dc-ch2-2-6-exist};
  linearity holds because the parallelism system $\dot V=A(t)V$, $A(t)=-\Gamma(\dot u,\cdot)(u)$,
  is linear, so solutions superpose and, by \cref{lem:dc-ch2-2-6-unique}, the endpoint map is
  additive and homogeneous. Injectivity follows from backward uniqueness of the same linear
  system (two fields sharing $V(b)$ share $V(a)$); an injective linear endomorphism of the
  finite-dimensional tangent space is an isomorphism, so $P_c$ is invertible, with inverse the
  parallel transport along the reversed curve.
\end{lemma}

\section{Riemannian connections}

\begin{definition}[compatible with the metric]
  \label{def:dc-ch2-3-1}
  \dcref{ch2:3.1}
  \lean{Riemannian.CompatibleGen}
  \leanok
  Let $M$ be a differentiable manifold with an affine connection $\nabla$ and a
  Riemannian metric $\langle\,,\,\rangle$. A connection is said to be \emph{compatible}
  with the metric $\langle\,,\,\rangle$ when, for any smooth curve $c$ and any pair
  of parallel vector fields $P$ and $P'$ along $c$, we have
  $\langle P,P'\rangle=\text{constant}$.
\end{definition}

\begin{proposition}[product rule characterization]
  \label{prop:dc-ch2-3-2}
  \dcref{ch2:3.2}
  \uses{def:dc-ch2-3-1, prop:dc-ch2-2-6, lem:dc-ch2-3-2-coord}
  \lean{Riemannian.compatibleGen_iff_productRuleGen}
  \leanok
  Let $M$ be a Riemannian manifold. A connection $\nabla$ on $M$ is compatible with
  a metric if and only if for any vector fields $V$ and $W$ along the differentiable
  curve $c:I\to M$ we have
  
  $$
  \frac{d}{dt}\langle V,W\rangle=\Big\langle\frac{DV}{dt},W\Big\rangle+\Big\langle V,\frac{DW}{dt}\Big\rangle,\quad t\in I.\qquad(3)
  $$
\end{proposition}
\begin{proof}
  \leanok
  Equation (3) obviously implies $\nabla$ is compatible. For the converse,
  choose an orthonormal basis $\{P_1(t_0),\dots,P_n(t_0)\}$ of $T_{c(t_0)}M$. Using
  \cref{prop:dc-ch2-2-6}, extend the $P_i(t_0)$ along $c$ by parallel translation; since
  $\nabla$ is compatible with the metric (\cref{def:dc-ch2-3-1}), $\{P_1(t),\dots,P_n(t)\}$
  is an orthonormal basis of $T_{c(t)}M$ for each $t$. Writing $V=\sum_i v^i P_i$,
  $W=\sum_i w^i P_i$ gives $\frac{DV}{dt}=\sum_i\frac{dv^i}{dt}P_i$,
  $\frac{DW}{dt}=\sum_i\frac{dw^i}{dt}P_i$, whence
  $\langle\frac{DV}{dt},W\rangle+\langle V,\frac{DW}{dt}\rangle
  =\sum_i(\frac{dv^i}{dt}w^i+\frac{dw^i}{dt}v^i)
  =\frac{d}{dt}\sum_i v^i w^i=\frac{d}{dt}\langle V,W\rangle$. $\blacksquare$
\end{proof}

\begin{corollary}[compatibility on vector fields]
  \label{cor:dc-ch2-3-3}
  \dcref{ch2:3.3}
  \lean{Riemannian.AffineConnection.IsMetricCompatible}
  \leanok
  \uses{prop:dc-ch2-3-2}
  A connection $\nabla$ on a Riemannian manifold $M$ is compatible with the metric
  if and only if
  
  $$
  X\langle Y,Z\rangle=\langle\nabla_X Y,Z\rangle+\langle Y,\nabla_X Z\rangle,\quad X,Y,Z\in\mathcal{X}(M).\qquad(4)
  $$
\end{corollary}
\begin{proof}
  Suppose $\nabla$ is compatible with the metric. Let $p\in M$ and
  $c:I\to M$ a curve with $c(t_0)=p$ and $\frac{dc}{dt}\big|_{t=t_0}=X(p)$. Then by
  \cref{prop:dc-ch2-3-2},
  $X(p)\langle Y,Z\rangle=\frac{d}{dt}\langle Y,Z\rangle\big|_{t=t_0}
  =\langle\nabla_{X(p)}Y,Z\rangle_p+\langle Y,\nabla_{X(p)}Z\rangle_p$. Since $p$ is
  arbitrary, (4) follows. The converse is obvious. $\blacksquare$
\end{proof}

\begin{lemma}[metric compatibility of $D/dt$, coordinate form]
  \label{lem:dc-ch2-3-2-coord}
  \dcref{ch2:3.2}
  \lean{Riemannian.hasDerivAt_chartMetricInner_along}
  \leanok
  \uses{lem:dc-ch2-2-2-coord, def:dc-ch2-3-7-christoffel}
  Read in a coordinate chart at $\alpha$, let $\langle a,b\rangle_y=\sum_{i,j}g_{ij}(y)\,a^i b^j$
  be the metric inner product on coordinate vectors and let $\frac{D}{dt}$ be the covariant
  derivative along a curve of \cref{lem:dc-ch2-2-2-coord}, $\frac{DV}{dt}=\dot V+\Gamma(\dot u,V)(u)$,
  built from the Levi-Civita chart Christoffel symbols. For a coordinate curve $u$ in the chart
  domain and coordinate vector fields $V,W$ along it,
  $$
  \frac{d}{dt}\langle V,W\rangle=\Big\langle\frac{DV}{dt},W\Big\rangle+\Big\langle V,\frac{DW}{dt}\Big\rangle.
  $$
  This is do Carmo's equation (3) (\cref{prop:dc-ch2-3-2}) for the Levi-Civita connection.
\end{lemma}
\begin{proof}
  \leanok
  \uses{def:dc-ch2-3-7-christoffel}
  Differentiate $t\mapsto\sum_{i,j}g_{ij}(u(t))\,V^i(t)\,W^j(t)$ by the product and chain rules:
  $\frac{d}{dt}\langle V,W\rangle=\sum_{i,j}\big(\sum_k\dot u^k\,\partial_k g_{ij}\big)V^iW^j
  +\sum_{i,j}g_{ij}\dot V^iW^j+\sum_{i,j}g_{ij}V^i\dot W^j$. Substitute the coordinate form of
  $\nabla g=0$, do Carmo's formula (10) $\partial_k g_{ij}=\sum_m(g_{mj}\Gamma^m_{ki}+g_{im}\Gamma^m_{kj})$
  (\cref{def:dc-ch2-3-7-christoffel}). Reindexing the two resulting Christoffel sums identifies them
  with $\langle\Gamma(\dot u,V),W\rangle$ and $\langle V,\Gamma(\dot u,W)\rangle$, and combining with
  the $\dot V,\dot W$ terms gives $\langle\dot V+\Gamma(\dot u,V),W\rangle+\langle V,\dot W+\Gamma(\dot u,W)\rangle
  =\langle\frac{DV}{dt},W\rangle+\langle V,\frac{DW}{dt}\rangle$.
\end{proof}

\begin{lemma}[parallel fields have constant inner product, coordinate form]
  \label{lem:dc-ch2-3-1-coord}
  \dcref{ch2:3.1}
  \lean{Riemannian.hasDerivAt_chartMetricInner_eq_zero_of_isParallelCoord}
  \leanok
  \uses{def:dc-ch2-2-5, def:dc-ch2-3-1, lem:dc-ch2-3-2-coord}
  Read in a coordinate chart at $\alpha$, if $V$ and $W$ are parallel (\cref{def:dc-ch2-2-5})
  coordinate vector fields along a curve $u$ in the chart domain, then $\langle V,W\rangle$ is
  constant along $u$: its derivative vanishes identically. Thus the Levi-Civita connection is
  compatible with the metric in do Carmo's sense (\cref{def:dc-ch2-3-1}).
\end{lemma}
\begin{proof}
  \leanok
  \uses{lem:dc-ch2-3-2-coord}
  Immediate from \cref{lem:dc-ch2-3-2-coord}: both covariant derivatives $\frac{DV}{dt}$ and
  $\frac{DW}{dt}$ vanish, so $\frac{d}{dt}\langle V,W\rangle
  =\langle 0,W\rangle+\langle V,0\rangle=0$.
\end{proof}

\begin{definition}[symmetric connection]
  \label{def:dc-ch2-3-4}
  \dcref{ch2:3.4}
  \lean{Riemannian.AffineConnection.IsSymmetric}
  \leanok
  \uses{def:dc-ch2-2-1}
  An affine connection $\nabla$ on a smooth manifold $M$ is said to be \emph{symmetric}
  when
  
  $$
  \nabla_X Y-\nabla_Y X=[X,Y]\quad\text{for all }X,Y\in\mathcal{X}(M).\qquad(5)
  $$
\end{definition}

\begin{remark}[Remark 3.5]
  \label{rem:dc-ch2-3-5}
  \dcref{ch2:3.5}
  In a coordinate system $(U,\mathbf{x})$, symmetry of $\nabla$ implies for all
  $i,j=1,\dots,n$,
  
  $$
  \nabla_{X_i}X_j-\nabla_{X_j}X_i=[X_i,X_j]=0,\quad X_i=\frac{\partial}{\partial x_i},\qquad(5')
  $$
  
  which justifies the terminology; (5') is equivalent to $\Gamma_{ij}^k=\Gamma_{ji}^k$.
\end{remark}

\begin{lemma}[Koszul formula]
  \label{lem:dc-ch2-3-6-koszul}
  \dcref{ch2:3.6}
  \lean{Riemannian.AffineConnection.koszul_formula}
  \leanok
  \uses{def:dc-ch2-3-4, cor:dc-ch2-3-3}
  Let $M$ be a Riemannian manifold with an affine connection $\nabla$ that is
  symmetric (\cref{def:dc-ch2-3-4}) and compatible with the metric in the sense of
  (4) (\cref{cor:dc-ch2-3-3}). Then for all $X,Y,Z\in\mathcal{X}(M)$,

  $$
  2\langle Z,\nabla_Y X\rangle=X\langle Y,Z\rangle+Y\langle Z,X\rangle-Z\langle X,Y\rangle-\langle[X,Z],Y\rangle-\langle[Y,Z],X\rangle-\langle[X,Y],Z\rangle.\qquad(9)
  $$
\end{lemma}
\begin{proof}
  \leanok
  \uses{def:dc-ch2-3-4, cor:dc-ch2-3-3}
  Applying (4) to the three cyclic orderings of $X,Y,Z$ gives (6), (7), (8).
  Adding (6) and (7), subtracting (8), replacing
  $\nabla_X Y-\nabla_Y X=[X,Y]$ and its two cyclic analogues by the symmetry (5) of
  $\nabla$ (\cref{def:dc-ch2-3-4}), and using the symmetry of $\langle\,,\,\rangle$
  to pair the remaining terms, yields (9).
\end{proof}

\begin{lemma}[tangent-vector extension]
  \label{lem:dc-ch2-tangent-extension}
  \lean{Riemannian.exists_smoothVectorField_eq}
  \leanok
  On a finite-dimensional, $\sigma$-compact, Hausdorff smooth manifold $M$, every
  tangent vector $v\in T_pM$ is the value at $p$ of some global smooth vector field
  $Z\in\mathcal{X}(M)$; that is, $Z(p)=v$.
\end{lemma}
\begin{proof}
  \leanok
  Let $t_x=\{v\}$ for $x=p$ and $t_x=T_xM$ otherwise; each $t_x$ is convex. Near
  $p$, the section $x\mapsto e^{-1}(x,c)$ that is constant in a bundle
  trivialization $e$ around $p$ (with $c$ the fibre coordinate of $v$) is smooth and
  equals $v$ at $p$; away from $p$ the zero section lies in $t_x$. Gluing these
  local sections with a partition of unity yields a global smooth field $Z$ with
  $Z(p)\in\{v\}$, i.e.\ $Z(p)=v$.
\end{proof}

\begin{lemma}[uniqueness of the Levi-Civita connection]
  \label{lem:dc-ch2-3-6-unique}
  \dcref{ch2:3.6}
  \lean{Riemannian.AffineConnection.leviCivita_unique'}
  \leanok
  \uses{lem:dc-ch2-3-6-koszul, lem:dc-ch2-tangent-extension}
  On a finite-dimensional, $\sigma$-compact, Hausdorff Riemannian manifold there is
  \emph{at most one} affine connection that is symmetric (\cref{def:dc-ch2-3-4})
  and compatible with the metric in the sense of (4) (\cref{cor:dc-ch2-3-3}).
\end{lemma}
\begin{proof}
  \leanok
  \uses{lem:dc-ch2-3-6-koszul, lem:dc-ch2-tangent-extension}
  If $\nabla_1,\nabla_2$ are both symmetric and metric-compatible then, by the
  Koszul formula (\cref{lem:dc-ch2-3-6-koszul}), $\langle Z,(\nabla_1)_YX\rangle=
  \langle Z,(\nabla_2)_YX\rangle$ for every $Z\in\mathcal{X}(M)$: both sides equal
  the metric-only right-hand side of (9). By \cref{lem:dc-ch2-tangent-extension}
  every tangent vector is realized as some $Z(p)$, so the two covariant derivatives
  have equal inner product against every tangent vector; non-degeneracy of the
  metric then forces $(\nabla_1)_YX=(\nabla_2)_YX$, hence $\nabla_1=\nabla_2$.
\end{proof}

\begin{lemma}[well-definedness of the Koszul functional]
  \label{lem:dc-ch2-3-6-welldef}
  \dcref{ch2:3.6}
  \lean{Riemannian.RiemannianMetric.koszulRHS,
    Riemannian.RiemannianMetric.koszulRHS_add_right,
    Riemannian.RiemannianMetric.koszulRHS_smul_right}
  \leanok
  \uses{lem:dc-ch2-3-6-koszul}
  Fix a Riemannian manifold $M$ and $X, Y \in \mathcal{X}(M)$. The metric-only
  right-hand side of the Koszul formula (9),
  \[
    F_{Y,X}(Z) = X\langle Y,Z\rangle + Y\langle Z,X\rangle - Z\langle X,Y\rangle
      - \langle[X,Z],Y\rangle - \langle[Y,Z],X\rangle - \langle[X,Y],Z\rangle,
  \]
  is $\mathcal{D}(M)$-linear in the test field $Z$: it is additive,
  $F_{Y,X}(Z_1+Z_2) = F_{Y,X}(Z_1) + F_{Y,X}(Z_2)$, and function-homogeneous,
  $F_{Y,X}(fZ) = f\,F_{Y,X}(Z)$ for every $f \in \mathcal{D}(M)$. Consequently
  $Z \mapsto \tfrac12 F_{Y,X}(Z)$ is a tensorial functional, and its metric Riesz
  dual is the candidate covariant derivative $\nabla_Y X$ for the existence half of
  \cref{thm:dc-ch2-3-6}.
\end{lemma}
\begin{proof}
  \leanok
  \uses{lem:dc-ch2-3-6-koszul}
  Additivity is immediate from additivity of the directional derivative, of the
  metric in each slot, and of the Lie bracket in the second slot. For homogeneity,
  the Leibniz rule $X(fh) = f\,X(h) + h\,X(f)$ and the bracket Leibniz rule
  $[X,fZ] = f\,[X,Z] + X(f)\,Z$ produce the extra terms $X(f)\langle Y,Z\rangle$
  and $-X(f)\langle Z,Y\rangle$ (together with the analogous $Y(f)$ pair), which
  cancel by symmetry of the metric $\langle Y,Z\rangle = \langle Z,Y\rangle$.
\end{proof}

\begin{lemma}[the Koszul-dual connection is an affine connection]
  \label{lem:dc-ch2-3-6-affine}
  \dcref{ch2:3.6}
  \lean{Riemannian.RiemannianMetric.koszulRHS_add_left,
    Riemannian.RiemannianMetric.koszulRHS_leibniz_left,
    Riemannian.RiemannianMetric.koszulRHS_add_middle,
    Riemannian.RiemannianMetric.koszulRHS_smul_middle}
  \leanok
  \uses{def:dc-ch2-2-1, lem:dc-ch2-3-6-welldef}
  Fix a Riemannian manifold $M$ and a test field $Z \in \mathcal{X}(M)$. As a
  function of its two connection arguments, the Koszul functional
  $F_{Y,X}(Z) = X\langle Y,Z\rangle + Y\langle Z,X\rangle - Z\langle X,Y\rangle
    - \langle[X,Z],Y\rangle - \langle[Y,Z],X\rangle - \langle[X,Y],Z\rangle$
  realizes the axioms of an affine connection (\cref{def:dc-ch2-2-1}) for the
  candidate covariant derivative $\nabla_Y X := \tfrac12 (F_{Y,X})^\sharp$:
  \begin{itemize}
    \item additivity in the differentiated field,
      $F_{Y,X_1+X_2} = F_{Y,X_1} + F_{Y,X_2}$;
    \item the Leibniz rule in the differentiated field,
      $F_{Y,fX} = f\,F_{Y,X} + 2\,Y(f)\,\langle X,Z\rangle$
      (i.e.\ $\nabla_Y(fX) = f\,\nabla_Y X + Y(f)\,X$);
    \item additivity in the direction field,
      $F_{Y_1+Y_2,X} = F_{Y_1,X} + F_{Y_2,X}$;
    \item $\mathcal{D}(M)$-homogeneity in the direction field,
      $F_{fY,X} = f\,F_{Y,X}$ (i.e.\ $\nabla_{fY}X = f\,\nabla_Y X$).
  \end{itemize}
\end{lemma}
\begin{proof}
  \leanok
  \uses{lem:dc-ch2-3-6-koszul}
  Each identity is obtained by substituting the modified field into the six terms of
  $F$ and applying additivity/homogeneity of the directional derivative, the
  bilinearity of the metric, and the corresponding bracket rule. Additivity in both
  slots is immediate. In the differentiated field, the field homogeneity
  $(fX)(h) = f\,X(h)$, the metric Leibniz rule and the bracket Leibniz rules
  $[fX,W] = f\,[X,W] - W(f)\,X$ produce four boundary terms; two cancel against each
  other and the remaining pair, $Y(f)\langle Z,X\rangle + Y(f)\langle X,Z\rangle$,
  combines to $2\,Y(f)\langle X,Z\rangle$ by symmetry of the metric. In the direction
  field all four boundary terms cancel in pairs by symmetry of the metric, leaving
  pure homogeneity. Together with the $Z$-linearity of
  \cref{lem:dc-ch2-3-6-welldef}, this is do Carmo's remark that $\nabla$ defined by
  (9) is a well-defined affine connection; only the smoothness of the Riesz-dual
  section remains for existence.
\end{proof}

\begin{lemma}[the Koszul-dual connection is Levi-Civita]
  \label{lem:dc-ch2-3-6-exist}
  \dcref{ch2:3.6}
  \lean{Riemannian.AffineConnection.isLeviCivita_of_koszulDual,
    Riemannian.koszulDual_isSymmetric,
    Riemannian.koszulDual_isMetricCompatible}
  \leanok
  \uses{def:dc-ch2-2-1, def:dc-ch2-3-4, lem:dc-ch2-3-6-koszul, lem:dc-ch2-3-6-welldef,
    lem:dc-ch2-tangent-extension}
  On a finite-dimensional, $\sigma$-compact, Hausdorff Riemannian manifold, suppose
  $\nabla$ is an affine connection (\cref{def:dc-ch2-2-1}) whose covariant derivative
  is the metric Riesz dual of the Koszul functional (9): $2\langle\nabla_Y X, Z\rangle
  = F_{Y,X}(Z)$ for every test field $Z$. Then $\nabla$ is a Levi-Civita connection
  for the metric — it is symmetric (\cref{def:dc-ch2-3-4}) and compatible with the
  metric (\cref{cor:dc-ch2-3-3}).
\end{lemma}
\begin{proof}
  \leanok
  \uses{lem:dc-ch2-3-6-koszul, lem:dc-ch2-tangent-extension}
  Two algebraic identities for the Koszul functional close the two conditions.
  Swapping its connection slots, $F_{Y,X}(Z)-F_{X,Y}(Z)=2\langle[X,Y],Z\rangle$: the
  metric-derivative terms cancel by symmetry of the metric and the bracket terms
  combine via $[Y,X]=-[X,Y]$. With the defining identity this reads
  $2\langle\nabla_X Y-\nabla_Y X, Z\rangle=2\langle[X,Y],Z\rangle$ for all $Z$, so by
  the tangent-vector extension (\cref{lem:dc-ch2-tangent-extension}) and
  non-degeneracy of the metric, $\nabla_X Y-\nabla_Y X=[X,Y]$; that is symmetry.
  Adding the two orderings that share the base field $X$,
  $F_{Y,X}(Z)+F_{Z,X}(Y)=2\,X\langle Y,Z\rangle$: the bracket terms cancel in pairs by
  antisymmetry. With the defining identity this is exactly
  $X\langle Y,Z\rangle=\langle\nabla_X Y, Z\rangle+\langle Y,\nabla_X Z\rangle$, which
  is metric-compatibility (\cref{cor:dc-ch2-3-3}).
\end{proof}

\begin{theorem}[Levi-Civita]
  \label{thm:dc-ch2-3-6}
  \dcref{ch2:3.6}
  % Lean: unconditional existence+uniqueness. The smoothness of the Koszul-Riesz
  % dual section (do Carmo's "easy to verify") is discharged via the chart-frame
  % musical-isomorphism smoothness engine, which is set up for a model space
  % carrying an inner product; the Lean statement therefore additionally assumes
  % the model E is an inner-product space and finrank E ≠ 0 (WLOG for a
  % finite-dimensional manifold).
  \lean{Riemannian.RiemannianMetric.exists_unique_isLeviCivita}
  \leanok
  \uses{def:dc-ch2-3-4, cor:dc-ch2-3-3, lem:dc-ch2-3-6-koszul, lem:dc-ch2-3-6-unique, lem:dc-ch2-3-6-welldef, lem:dc-ch2-3-6-exist}
  Given a Riemannian manifold $M$, there exists a unique affine connection $\nabla$
  on $M$ satisfying the conditions:
  - a) $\nabla$ is symmetric.
  - b) $\nabla$ is compatible with the Riemannian metric.
\end{theorem}
\begin{proof}
  \leanok
  \uses{lem:dc-ch2-3-6-unique, lem:dc-ch2-3-6-exist}
  Suppose such a $\nabla$ exists. By \cref{cor:dc-ch2-3-3},
  
  $$
  X\langle Y,Z\rangle=\langle\nabla_X Y,Z\rangle+\langle Y,\nabla_X Z\rangle,\qquad(6)
  $$
  
  $$
  Y\langle Z,X\rangle=\langle\nabla_Y Z,X\rangle+\langle Z,\nabla_Y X\rangle,\qquad(7)
  $$
  
  $$
  Z\langle X,Y\rangle=\langle\nabla_Z X,Y\rangle+\langle X,\nabla_Z Y\rangle.\qquad(8)
  $$
  
  Adding (6) and (7), subtracting (8), and using the symmetry of $\nabla$
  (\cref{def:dc-ch2-3-4}) gives
  
  $$
  X\langle Y,Z\rangle+Y\langle Z,X\rangle-Z\langle X,Y\rangle=\langle[X,Z],Y\rangle+\langle[Y,Z],X\rangle+\langle[X,Y],Z\rangle+2\langle Z,\nabla_Y X\rangle.
  $$
  
  Therefore (the \emph{Koszul formula})
  
  $$
  \langle Z,\nabla_Y X\rangle=\tfrac{1}{2}\big\{X\langle Y,Z\rangle+Y\langle Z,X\rangle-Z\langle X,Y\rangle-\langle[X,Z],Y\rangle-\langle[Y,Z],X\rangle-\langle[X,Y],Z\rangle\big\}.\qquad(9)
  $$
  
  Expression (9) shows $\nabla$ is uniquely determined from the metric, so if it
  exists it is unique. For existence, define $\nabla$ by (9); it is easy to verify
  that $\nabla$ is well-defined and satisfies the desired conditions. $\blacksquare$
\end{proof}

\begin{definition}[Christoffel symbols]
  \label{def:dc-ch2-3-7-christoffel}
  \dcref{ch2:3.7}
  \lean{Riemannian.chartChristoffel}
  \leanok
  \uses{def:dc-ch1-2-1}
  In a coordinate neighborhood $U$ with metric coefficients
  $g_{ij}=\langle X_i,X_j\rangle$ and inverse matrix $(g^{km})$, the
  \emph{Christoffel symbols} (of the second kind) are the functions on $U$ given
  by the classical formula

  $$
  \Gamma_{ij}^m=\tfrac{1}{2}\sum_k\Big\{\frac{\partial}{\partial x_i}g_{jk}+\frac{\partial}{\partial x_j}g_{ki}-\frac{\partial}{\partial x_k}g_{ij}\Big\}g^{km}.\qquad(10)
  $$
\end{definition}

\begin{remark}[the Levi-Civita connection and its coefficients]
  \label{rem:dc-ch2-3-7}
  \dcref{ch2:3.7}
  \uses{thm:dc-ch2-3-6,def:dc-ch2-3-7-christoffel}
  The connection given by \cref{thm:dc-ch2-3-6} will be referred to, from now on, as the
  \emph{Levi-Civita} (or \emph{Riemannian}) \emph{connection} on $M$. The functions $\Gamma_{ij}^k$
  defined on $U$ by $\nabla_{X_i}X_j=\sum_k\Gamma_{ij}^k X_k$ are the
  \emph{coefficients of the connection} $\nabla$ on $U$, and by uniqueness in
  \cref{thm:dc-ch2-3-6} they agree with the Christoffel symbols of
  \cref{def:dc-ch2-3-7-christoffel}: from (9) it follows that

  $$
  \sum_\ell\Gamma_{ij}^\ell g_{\ell k}=\tfrac{1}{2}\Big\{\frac{\partial}{\partial x_i}g_{jk}+\frac{\partial}{\partial x_j}g_{ki}-\frac{\partial}{\partial x_k}g_{ij}\Big\},
  $$

  which is equivalent to (10). For Euclidean space $\mathbb{R}^n$ one has
  $\Gamma_{ij}^k=0$. In terms of the Christoffel symbols, the covariant derivative
  has the classical expression

  $$
  \frac{DV}{dt}=\sum_k\Big\{\frac{dv^k}{dt}+\sum_{i,j}\Gamma_{ij}^k v^j\frac{dx_i}{dt}\Big\}X_k,
  $$

  which follows from (1). In Euclidean spaces the covariant derivative coincides
  with the usual derivative.
\end{remark}
